\documentclass[10pt,a4paper]{amsart}
\usepackage{amsmath,amssymb,amsthm,graphicx}
% Retained minimal preamble for the Stacks Project chapter varieties.tex.
% Verbatim excerpt of the chapter preamble: Xy-pic setup (native preamble lines 14-18), the native theorem
% declarations with the mathematical macros (native preamble lines 154-217), and the project own no-op
% definitions of the reference/slogan/history display environments (native wrapper.tex lines 1-3) that the
% counted statement uses for its slogan. Copied mechanically; nothing added, removed or redefined.
% For commutative diagrams we use Xy-pic
\usepackage[all]{xy}

% We use 2cell for 2-commutative diagrams.
\xyoption{2cell}
% Theorem environments.
%
\theoremstyle{plain}
\newtheorem{theorem}[subsection]{Theorem}
\newtheorem{proposition}[subsection]{Proposition}
\newtheorem{lemma}[subsection]{Lemma}

\theoremstyle{definition}
\newtheorem{definition}[subsection]{Definition}
\newtheorem{example}[subsection]{Example}
\newtheorem{exercise}[subsection]{Exercise}
\newtheorem{situation}[subsection]{Situation}

\theoremstyle{remark}
\newtheorem{remark}[subsection]{Remark}
\newtheorem{remarks}[subsection]{Remarks}

\numberwithin{equation}{subsection}

% Macros
%
\def\lim{\mathop{\mathrm{lim}}\nolimits}
\def\colim{\mathop{\mathrm{colim}}\nolimits}
\def\Spec{\mathop{\mathrm{Spec}}}
\def\Hom{\mathop{\mathrm{Hom}}\nolimits}
\def\Ext{\mathop{\mathrm{Ext}}\nolimits}
\def\SheafHom{\mathop{\mathcal{H}\!\mathit{om}}\nolimits}
\def\SheafExt{\mathop{\mathcal{E}\!\mathit{xt}}\nolimits}
\def\Sch{\mathit{Sch}}
\def\Mor{\mathop{\mathrm{Mor}}\nolimits}
\def\Ob{\mathop{\mathrm{Ob}}\nolimits}
\def\Sh{\mathop{\mathit{Sh}}\nolimits}
\def\NL{\mathop{N\!L}\nolimits}
\def\CH{\mathop{\mathrm{CH}}\nolimits}
\def\proetale{{pro\text{-}\acute{e}tale}}
\def\etale{{\acute{e}tale}}
\def\QCoh{\mathit{QCoh}}
\def\Ker{\mathop{\mathrm{Ker}}}
\def\Im{\mathop{\mathrm{Im}}}
\def\Coker{\mathop{\mathrm{Coker}}}
\def\Coim{\mathop{\mathrm{Coim}}}

% Boxtimes
%
\DeclareMathSymbol{\boxtimes}{\mathbin}{AMSa}{"02}

%
% Macros for moduli stacks/spaces
%
\def\QCohstack{\mathcal{QC}\!\mathit{oh}}
\def\Cohstack{\mathcal{C}\!\mathit{oh}}
\def\Spacesstack{\mathcal{S}\!\mathit{paces}}
\def\Quotfunctor{\mathrm{Quot}}
\def\Hilbfunctor{\mathrm{Hilb}}
\def\Curvesstack{\mathcal{C}\!\mathit{urves}}
\def\Polarizedstack{\mathcal{P}\!\mathit{olarized}}
\def\Complexesstack{\mathcal{C}\!\mathit{omplexes}}
% \Pic is the operator that assigns to X its picard group, usage \Pic(X)
% \Picardstack_{X/B} denotes the Picard stack of X over B
% \Picardfunctor_{X/B} denotes the Picard functor of X over B
\def\Pic{\mathop{\mathrm{Pic}}\nolimits}
\def\Picardstack{\mathcal{P}\!\mathit{ic}}
\def\Picardfunctor{\mathrm{Pic}}
\def\Deformationcategory{\mathcal{D}\!\mathit{ef}}
\newenvironment{reference}{}{}
\newenvironment{slogan}{}{}
\newenvironment{history}{}{}
\usepackage{cleveref}
\crefname{originallemma}{original lemma}{original lemmas}
\Crefname{originallemma}{Original lemma}{Original lemmas}
\makeatletter
\newlabel{algebra-lemma-geometrically-irreducible}{{037K}{source}{}{sourcepointer-1}{}}
\newlabel{algebra-lemma-geometrically-irreducible@cref}{{[originallemma][0][]037K}{[1][0][]source}{}{}{}}
\makeatother
\crefname{originallemma}{original lemma}{original lemmas}
\Crefname{originallemma}{Original lemma}{Original lemmas}
\makeatletter
\newlabel{lemma-separably-closed-field-irreducible-components}{{038H}{source}{}{sourcepointer-2}{}}
\newlabel{lemma-separably-closed-field-irreducible-components@cref}{{[originallemma][0][]038H}{[1][0][]source}{}{}{}}
\makeatother
\crefname{originallemma}{original lemma}{original lemmas}
\Crefname{originallemma}{Original lemma}{Original lemmas}
\makeatletter
\newlabel{properties-lemma-characterize-irreducible}{{01OM}{source}{}{sourcepointer-3}{}}
\newlabel{properties-lemma-characterize-irreducible@cref}{{[originallemma][0][]01OM}{[1][0][]source}{}{}{}}
\makeatother
\title{Geometric irreducibility can be tested over a separable algebraic closure}
\author{The Stacks Project Authors}
\date{stacks-project:varieties.tex:local-snapshot-20261011; GNU FDL 1.2 or later}
\begin{document}
\maketitle
\noindent\textit{Editorial note.} All delivered mathematical text is the author's own, copied from the pinned Stacks Project chapter file varieties.tex; no statement or proof step was rewritten, reordered or removed.\par\medskip
\noindent\textit{Editorial note.} The statement carries the chapter's slogan environment, rendered by the no-op display environment defined at lines 75-77 of the retained preamble exactly as the author wrote it.\par\medskip
\noindent\textit{Editorial note.} The counted spans carry 3 navigation stubs and no bibliography citation.\par\medskip
\noindent\textit{Editorial note.} The preamble is the chapter's retained minimal preamble (sources/stacks-varieties/retained-preamble.tex lines 1-77), a verbatim excerpt of the native preamble.\par\medskip
\noindent\textit{Editorial note.} The label lemma-characterize-geometrically-irreducible is defined only in varieties.tex, so no cross-chapter restatement had to be weighed for this card.\par\medskip
\noindent\textit{Allowed standard prerequisites (editorial).}\begin{itemize}
\item Irreducible components of a base change to a separable algebraic closure. Same chapter, Lemma lemma-separably-closed-field-irreducible-components, invoked at proof lines 1563-1564. Supplies that irreducibility is preserved by extending the base field to a larger separably closed field.
\item Criterion for irreducibility of a scheme via an affine open covering. Properties, Lemma properties-lemma-characterize-irreducible, invoked at proof lines 1575-1578. Supplies that a reducible member or an empty intersection of two members makes X reducible.
\item Geometric irreducibility of a finite type algebra tested over a finite separable extension. Algebra, Lemma algebra-lemma-geometrically-irreducible, invoked at proof lines 1583-1586. Supplies the finite separable extension over which an affine member becomes reducible.
\end{itemize}\medskip
\ifdefined\hypertarget\hypertarget{sourcepointer-1}{}\fi
\noindent\textit{Source reference.} Original lemma 037K: sources/stacks-chow/tags.txt, source lines 4203--4203. Pointer outside the counted span, invoked once for context; the load-bearing argument of the counted statement is the author's proof reproduced in full in the proof span.\par\medskip
\ifdefined\hypertarget\hypertarget{sourcepointer-2}{}\fi
\noindent\textit{Source reference.} Original lemma 038H: sources/stacks-varieties/pinned-original-source.tex, source lines 1524--1524. Pointer outside the counted span, invoked once for context; the load-bearing argument of the counted statement is the author's proof reproduced in full in the proof span.\par\medskip
\ifdefined\hypertarget\hypertarget{sourcepointer-3}{}\fi
\noindent\textit{Source reference.} Original lemma 01OM: sources/stacks-chow/tags.txt, source lines 2199--2199. Pointer outside the counted span, invoked once for context; the load-bearing argument of the counted statement is the author's proof reproduced in full in the proof span.\par\medskip
\begin{lemma}
\label{lemma-characterize-geometrically-irreducible}
\begin{slogan}
Geometric irreducibility can be tested over a separable algebraic
closure of the base field.
\end{slogan}
Let $k$ be a field. Let $X$ be a scheme over $k$.
The following are equivalent:
\begin{enumerate}
\item $X$ is geometrically irreducible over $k$,
\item for every finite separable field extension $k'/k$
the scheme $X_{k'}$ is irreducible, and
\item $X_{\overline{k}}$ is irreducible, where $k \subset \overline{k}$
is a separable algebraic closure of $k$.
\end{enumerate}
\end{lemma}

\begin{proof}
Assume $X_{\overline{k}}$ is irreducible, i.e., assume (3).
Let $k'/k$ be a field extension.
There exists a field extension $\overline{k}'/\overline{k}$
such that $k'$ embeds into $\overline{k}'$ as an extension of $k$.
By Lemma \ref{lemma-separably-closed-field-irreducible-components}
we see that $X_{\overline{k}'}$ is irreducible.
Since $X_{\overline{k}'} \to X_{k'}$ is surjective we conclude
that $X_{k'}$ is irreducible. Hence (1) holds.

\medskip\noindent
Let $k \subset \overline{k}$ be a separable algebraic closure of $k$.
Assume not (3), i.e., assume $X_{\overline{k}}$ is reducible.
Our goal is to show that also $X_{k'}$ is
reducible for some finite subextension
$\overline{k}/k'/k$.
Let $X = \bigcup_{i \in I} U_i$ be an affine open covering
with $U_i$ not empty. If for some $i$ the scheme
$U_i$ is reducible, or if for some pair $i \not = j$ the
intersection $U_i \cap U_j$ is empty, then $X$ is reducible
(Properties, Lemma \ref{properties-lemma-characterize-irreducible})
and we are done.
In particular we may assume that
$U_{i, \overline{k}} \cap U_{j, \overline{k}}$ for all $i, j \in I$
is nonempty and we conclude that $U_{i, \overline{k}}$ has
to be reducible for some $i$. According to
Algebra, Lemma \ref{algebra-lemma-geometrically-irreducible}
this means that $U_{i, k'}$ is reducible for some
finite separable field extension $k'/k$.
Hence also $X_{k'}$ is reducible. Thus we see that
(2) implies (3).

\medskip\noindent
The implication (1) $\Rightarrow$ (2) is immediate.
This proves the lemma.
\end{proof}

\end{document}
